\chapter{Ретроспектива происхождения носителя \texorpdfstring{\(K\)}{K}}

После серии отрицательных gate-ов необходимо проверить, не были ли ограничения
созданы слишком ранней заморозкой геометрии. Для этого восстановим цепочку
от исходного zero-prompt и документа \texttt{TOE.pdf} до рабочего
\begin{equation}
 K=\mathbb{RP}^3\times S^1.
\end{equation}

\section{Что действительно делает zero-prompt}

Исходный файл
\path{RPFT-main/ai-promts/First-principles-00.md} содержит четыре шага.

Во-первых, он требует компактное односвязное трёхмерное group manifold
единичных кватернионов. Внутри уже выбранного класса это корректно даёт
\begin{equation}
 SU(2)\simeq S^3.
\end{equation}
Но трёхмерность, group-manifold form и quaternionic symmetry являются
входами prompt, а не результатом минимизации.

Во-вторых, prompt объявляет центр \(\{\pm1\}\) физически неразличимым и
требует antipodal quotient. После этого действительно
\begin{equation}
 S^3/\mathbb Z_2\simeq SO(3)\simeq\mathbb{RP}^3.
\end{equation}
Real structure \(J\) сама по себе такого quotient base points не доказывает;
необходим отдельный axiom о тривиальном действии центра на всех физических
состояниях и наблюдаемых.

В-третьих, из конечности KMS-системы prompt требует периодичность modular
flow и получает \(S^1\). Этот переход неверен в общем случае. Группа
Томиты--Такесаки имеет real parameter \(t\in\mathbb R\); конечномерная
density matrix может иметь несоизмеримые modular frequencies и не иметь
общего ненулевого периода.

Наконец, четвёртый пункт прямо предписывает product architecture. Поэтому
\begin{equation}
 \mathbb{RP}^3\times S^1
\end{equation}
является корректным ответом на поставленный prompt, но direct product уже
содержится в условии. Не сравниваются warped metrics и нетривиальный
\(U(1)\)-bundle class, существующий из-за
\begin{equation}
 H^2(\mathbb{RP}^3;\mathbb Z)\cong\mathbb Z_2.
\end{equation}

\begin{theorem}[Zero-prompt status]
Zero-prompt условно реконструирует \(K\), но не выводит его как уникальный
vacuum. Он заранее фиксирует quaternionic spatial class, central quotient,
periodic circle и product architecture; radius и variational stability не
определены.
\end{theorem}

\section{Что на самом деле постулируется в TOE}

Первичный документ \texttt{TOE.pdf} начинается не с компактного
\(K=\mathbb{RP}^3\times S^1\), а с compact self-adjoint correlation
operator \(\widehat C\), ядро которого определено на паре точек
\begin{equation}
 (x,y)\in M\times M.
\end{equation}
Здесь \(M\times M\) --- domain двухточечного kernel, а не утверждение, что
физическое пространство является произведением двух manifold factors.
Постулируются
\begin{align}
 \widehat C&=e^{-\sigma^2\Delta},
 \\
 C_\sigma(x,y)
 &=\frac1{(4\pi\sigma^2)^2}
 \exp\left(-\frac{|x-y|^2}{4\sigma^2}\right),
\end{align}
а metric восстанавливается из logarithmic Hessian kernel:
\begin{equation}
 g_{\mu\nu}(x)
 =-\lim_{y\to x}4\sigma^2
 \frac{\partial^2}{\partial x^\mu\partial y^\nu}
 \log C(x,y).
\end{equation}
Таким образом, TOE задаёт правило emergent local geometry, но не выбирает
топологию \(M\).

Standard Model sector в TOE использует другой объект: finite two-point
space \(F\) с
\begin{equation}
 A_F=\mathbb C\oplus\mathbb H\oplus M_3(\mathbb C),
 \qquad
 A=C^\infty(M)\otimes A_F.
\end{equation}
Это finite noncommutative factor, а не \(\mathbb{RP}^3\times S^1\).

Для трёх поколений TOE затем вводит ещё один, логически отдельный compact
internal manifold \(M_{\rm int}\) и заявляет
\begin{equation}
 N_{\rm gen}=\operatorname{Ind}(D)
 =\frac12|\chi(M_{\rm int})|,
 \qquad
 \chi(M_{\rm int})=\pm6.
\end{equation}
Но конкретное \(M_{\rm int}\) не строится. Поздний \texttt{TOE5.pdf}
прямо называет это Topological Realization Problem и заменяет manifold
лишь conjectural discrete multi-sheet space.

\begin{proposition}[Различение трёх пространств TOE]
В TOE нельзя отождествлять:
\begin{enumerate}
 \item \(M\times M\) --- область двухточечного correlation kernel;
 \item \(F\) --- finite NCG space Standard Model;
 \item \(M_{\rm int}\) --- гипотетический generation-index carrier.
\end{enumerate}
Ни один из этих объектов в первичном документе не равен рабочему
\(K=\mathbb{RP}^3\times S^1\).
\end{proposition}

\section{Где появляется \texorpdfstring{\(S^3\times S^1\)}{S3 x S1}}

Геометрия \(S^3\times S^1\) появляется в позднем
\path{toe_ugsm_common_shadow_bridge.tex} как минимальный контролируемый
heat-trace experiment. Она выбирается для проверки, может ли один и тот же
Gaussian/heat operator воспроизвести UGSM coefficients \(4\pi^3\),
\(\pi^2\) и linear \(\pi\)-proxy. Это диагностический background,
унаследованный от UGSM-стороны, а не решение equation of motion
\(\delta S[\widehat C]=0\).

Переход
\begin{equation}
 S^3\times S^1
 \longrightarrow
 \mathbb{RP}^3\times S^1
\end{equation}
добавляется позднее через RPFT/S2T central quotient и restricted carrier
criteria. Следовательно, фактическая genealogy имеет вид
\begin{equation}
 \boxed{
 \begin{gathered}
 \text{TOE: }\widehat C\text{ на общем }M\times M
 \\
 \downarrow\ \text{diagnostic choice}
 \\
 S^3\times S^1
 \downarrow\ \mathbb Z_2\text{-quotient и product assumptions}
 \\
 K=\mathbb{RP}^3\times S^1.
 \end{gathered}}
\end{equation}

\begin{theorem}[Carrier genealogy verdict]
Непрерывного вывода \(\text{TOE}\Rightarrow K\) в корпусе нет. Рабочий
\(K\) является осмысленным условным носителем, собранным из TOE heat-kernel
language, UGSM test geometry и RPFT/S2T quotient assumptions. Поэтому
отрицательное закрытие минимального fixed-\(K\) parent action не является
опровержением исходной correlation-operator программы TOE.
\end{theorem}

\section{Следующее направление исследования}

Ретроспектива меняет правильный следующий вопрос. Не следует немедленно
добавлять к фиксированному \(K\) ещё один сектор. Нужно вернуться к
первичному TOE-принципу и варьировать одновременно correlation operator и
carrier class:
\begin{equation}
 (M,\widehat C,\mathfrak s,\rho)
 \longmapsto
 \mathcal S_{\rm corr}[M,\widehat C;\mathfrak s,\rho].
\end{equation}
Только если один нормированный functional имеет уникальный устойчивый
minimum на \(\mathbb{RP}^3\times S^1\), слово ``выведенный носитель'' станет
допустимым. В противном случае следует сравнить surviving geometries, а не
принуждать их к уже известному \(K\).

\begin{protocol}[Новая конструктивная ветвь]
Следующий том может быть переоткрыт как variational carrier classification:
\begin{enumerate}
 \item объявить comparison class compact spin four-geometries и bundle
 sectors без вставки \(K\);
 \item определить один Gaussian-correlation spectral functional;
 \item вычислить его local coefficients и nonlocal determinant differences;
 \item экстремизировать radii, quotient и bundle classes;
 \item проверить global minimum и Hessian до observed-data readout.
\end{enumerate}
Это новая архитектурная ветвь, но она ближе к исходному TOE, чем дальнейшее
латание уже закрытого minimal fixed-\(K\) action.
\end{protocol}

Машинный ledger сохранён в
\path{s2t_v4_zero_prompt_toe_carrier_trace_gate_results.json}.